% Въ лѣто Господне ҂вк҃ѕ.
%
% Господи Іисусе Христе, Сыне Божій, помилуй меня грѣшнаго.
%
%
% Источникъ: https://azbyka.ru/molitvoslov/molitvy-na-son-grjadushhim-cerkovnoslavjanskim-shriftom.html

\documentclass[20pt]{extarticle}

\usepackage[a5paper, portrait, margin=.7cm, top=.3cm, bottom=.3cm]{geometry}

%\usepackage{pgfpages}
%\pgfpagesuselayout{2 on 1}[a4paper, landscape]

\usepackage{fontspec}
\usepackage{xkeyval}
\usepackage{microtype}

\usepackage{polyglossia}
\setmainlanguage{churchslavonic}

\usepackage{churchslavonic}
\usepackage{lettrine}
\usepackage{graphicx}  % for \scalebox

\usepackage{hyperref}
\hypersetup{
	pdftitle={Молитвы на сон грядущим},
	pdfauthor={Православная Церковь}
}


\makeatletter

\def\@texttop{\vfill}
\def\@textbottom{\vfill}

\makeatother


\newcommand{\cuL}[1]{\scalebox{1.5}{\cuKinovarColor#1}}
\newcommand{\cul}[1]{\mbox{\cuKinovarColor#1}}

\newfontfamily\noto{Noto Sans Symbols 2}
\newfontfamily\churchslavonicfont[Script=Cyrillic,Ligatures=TeX,HyphenChar=_]{Acathist}
%\newfontfamily\cyrillicfontsf{DejaVu Sans}

%\hyphenpenalty=10000
\setlength\emergencystretch{4em}
\def\hyph{{\_}{}{}\penalty-10000}

%\linespread{1.2}
\setlength{\parindent}{0cm}
\setlength{\parskip}{12pt}

\newcommand{\cuk}[1]{\cuKinovar{\normalsize #1}}
\newcommand{\sub}[1]{\vspace{4pt}{\centering\cuk{#1}\par}}


\begin{document}

\pagestyle{empty}
\large

\vspace*{-4cm}
\begin{center}
\cuKinovarColor{\Large\noto 🙚}
\begin{minipage}{.73\textwidth}\centering
	% TODO
	Моли́твы на со̀н грѧдꙋ́щимъ꙳  
\end{minipage}
{\Large\noto 🙘}
\end{center}
\vspace{-.7cm}

\cuL Во и҆́мѧ ѻ҆ц҃а̀, и҆ сн҃а, и҆ ст҃а́гѡ дх҃а. \cul{А҆}ми́нь.

\cuL Гдⷭ҇и і҆и҃се хрⷭ҇тѐ сн҃е бж҃їй, мл҃твъ ра́ди пречⷭ҇тыѧ твоеѧ̀ мт҃ре, прпⷣбныхъ и҆ бг҃оно́сныхъ ѻ҆тє́цъ на́шихъ, и҆ всѣ́хъ ст҃ы́хъ, поми́лꙋй на́съ. \cul{А҆}ми́нь.

\cuL С\hspace{-2pt}ла́ва тебѣ̀ бж҃е на́шъ, сла́ва тебѣ̀.

\cuL Цр҃ю̀ нбⷭ҇ный, ѹ҆тѣ́шителю, дш҃е и҆́стины, и҆́же вездѣ̀ сы́й, и҆ всѧ̑ и҆сполнѧ́ѧй, сокро́вище бл҃ги́хъ, и҆ жи́зни пода́телю, прїидѝ и҆ всели́сѧ въ ны̀, и҆ ѡ҆чи́сти ны̀ ѿ всѧ́кїѧ скве́рны, и҆ сп҃сѝ, бл҃же, дꙋ́шы на́шѧ.

\vspace{-8pt}
\cuL Ст҃ы́й бж҃е, ст҃ы́й крѣ́пкїй, ст҃ы́й безсме́ртный, поми́лꙋй на́съ. \cuk{Три́жды.}

\newpage
\cuL С\hspace{-2pt}ла́ва ѻ҆ц҃ꙋ̀, и҆ сн҃ꙋ, и҆ ст҃о́мꙋ дх҃ꙋ, и҆ ны́нѣ и҆ при́снѡ, и҆ во вѣ́ки вѣкѡ́въ. \cul{А҆}ми́нь.

\cuL Прест҃а́ѧ трⷪ҇це, поми́лꙋй на́съ: гдⷭ҇и, ѡ҆чи́сти грѣхѝ на́шѧ: влⷣко, простѝ беззакѡ́нїѧ на̑ша: ст҃ы́й, посѣтѝ и҆ и҆сцѣлѝ не́мѡщи на́шѧ, и҆́мене твоегѡ̀ ра́ди.

\vspace{-8pt}
\cuL Гдⷭ҇и, поми́лꙋй. \cuk{Три́жды.}

\cuL С\hspace{-2pt}ла́ва ѻ҆ц҃ꙋ̀, и҆ сн҃ꙋ, и҆ ст҃о́мꙋ дх҃ꙋ, и҆ ны́нѣ и҆ при́снѡ, и҆ во вѣ́ки вѣкѡ́въ. \cul{А҆}ми́нь.

\vspace{-14pt}
\cuL{Ѻ҆́}ч҃е на́шъ, и҆́же є҆сѝ на нб҃сѣ́хъ, да ст҃и́тсѧ и҆́мѧ твоѐ, да прїи́детъ црⷭ҇твїе твоѐ, да бꙋ́детъ во́лѧ твоѧ̀, ꙗ҆́кѡ на нб҃сѝ и҆ на землѝ. Хлѣ́бъ на́шъ насꙋ́щный да́ждь на́мъ дне́сь: и҆ ѡ҆ста́ви на́мъ до́лги на́шѧ, ꙗ҆́коже и҆ мы̀ ѡ҆ставлѧ́емъ должникѡ́мъ на́шымъ: и҆ не введѝ на́съ во и҆скꙋше́нїе, но и҆зба́ви на́съ ѿ лꙋка́вагѡ.

\newpage
\sub{Тропарѝ, гла́съ ѕ҃:}

Поми́лꙋй на́съ гдⷭ҇и, поми́лꙋй на́съ: всѧ́кагѡ бо ѿвѣ́та недоꙋмѣ́юще, сїю̀ тѝ мл҃твꙋ ꙗ҆́кѡ влⷣцѣ грѣ́шнїи прино́симъ: поми́лꙋй на́съ.

Сла́ва ѻ҆ц҃ꙋ̀, и҆ сн҃ꙋ, и҆ ст҃о́мꙋ дх҃ꙋ.

Гдⷭ҇и, поми́лꙋй на́съ, на тѧ́ бо ѹ҆пова́хомъ: не прогнѣ́вайсѧ на ны̀ ѕѣлѡ̀, нижѐ помѧнѝ беззако́нїй на́шихъ, но при́зри и҆ ны́нѣ ꙗ҆́кѡ бл҃гоꙋтро́бенъ, и҆ и҆зба́ви ны̀ ѿ вра̑гъ на́шихъ: ты́ бо є҆сѝ бг҃ъ на́шъ, и҆ мы̀ лю́дїе твоѝ, всѝ дѣла̀ рꙋкꙋ̀ твоє́ю, и҆ и҆́мѧ твоѐ призыва́емъ.

И҆ ны́нѣ и҆ при́снѡ, и҆ во вѣ́ки вѣкѡ́въ. А҆ми́нь.

Млⷭ҇рдїѧ двє́ри ѿве́рзи на́мъ, бл҃гослове́ннаѧ бцⷣе: надѣ́ющїисѧ на тѧ̀ да не поги́бнемъ, но да и҆зба́вимсѧ тобо́ю ѿ бѣ́дъ: ты́ бо є҆сѝ сп҃се́нїе ро́да хрⷭ҇тїа́нскагѡ.

\cuL Гдⷭ҇и, поми́лꙋй. \cuk{в҃i.}

\sub{Мл҃тва №, ст҃а́гѡ мака́рїа вели́кагѡ, къ бг҃ꙋ ѻ҆ц҃ꙋ̀:}
Бж҃е вѣ́чный и҆ цр҃ю̀ всѧ́кагѡ созда́нїѧ, сподо́бивый мѧ̀ да́же въ ча́съ се́й доспѣ́ти, прости́ ми грѣхѝ, ꙗ҆̀же сотвори́хъ въ се́й де́нь, дѣ́ломъ, сло́вомъ и҆ помышле́нїемъ: и҆ ѡ҆чи́сти, гдⷭ҇и, смире́ннꙋю мою̀ дꙋ́шꙋ ѿ всѧ́кїѧ скве́рны пло́ти и҆ дꙋ́ха, и҆ да́ждь мѝ, гдⷭ҇и, въ нощѝ се́й со́нъ прейтѝ въ ми́рѣ: да воста́въ ѿ смире́ннагѡ мѝ ло́жа, бл҃гоꙋгождꙋ̀ прест҃о́мꙋ и҆́мени твоемꙋ̀ во всѧ̑ дни̑ живота̀ моегѡ̀, и҆ поперꙋ̀ борю́щыѧ мѧ̀ врагѝ, плотскі̑ѧ и҆ безплѡ́тныѧ: и҆ и҆зба́ви мѧ̀, гдⷭ҇и, ѿ помышле́нїй сꙋ́етныхъ, ѡ҆сквернѧ́ющихъ мѧ̀, и҆ по́хотей лꙋка́выхъ. Ꙗ҆́кѡ твоѐ є҆́сть црⷭ҇тво, и҆ си́ла, и҆ сла́ва, ѻ҆ц҃а̀, и҆ сн҃а, и҆ ст҃а́гѡ дх҃а, ны́нѣ и҆ при́снѡ, и҆ во вѣ́ки вѣкѡ́въ. А҆ми́нь.

Мл҃тва в҃, ст҃а́гѡ а҆нтїо́ха,
ко гдⷭ҇ꙋ на́шемꙋ і҆и҃сꙋ хрⷭ҇тꙋ̀:
Вседержи́телю, сло́во ѻ҆́ч҃ее, са́мъ соверше́нъ сы́й, і҆и҃се хрⷭ҇тѐ, мно́гагѡ ра́ди млⷭ҇рдїѧ твоегѡ̀ никогда́же ѿлꙋча́йсѧ менє̀ раба̀ твоегѡ̀, но всегда̀ во мнѣ̀ почива́й. І҆и҃се, до́брый па́стырю твои́хъ ѻ҆ве́цъ, не преда́ждь менє̀ крамолѣ̀ ѕмїи́нѣ, и҆ жела́нїю сатанинꙋ̀ не ѡ҆ста́ви менє̀, ꙗ҆́кѡ сѣ́мѧ тлѝ в нем є҆́сть[*]. Ты̀ ѹ҆́бѡ гдⷭ҇и бж҃е покланѧ́емый, цр҃ю̀ ст҃ы́й і҆и҃се хрⷭ҇тѐ, спѧ́ща мѧ̀ сохранѝ немерца́ющимъ свѣ́томъ, дх҃омъ твои́мъ ст҃ы́мъ, и҆́мже ѡ҆ст҃и́лъ є҆сѝ твоѧ̑ ѹ҆ч҃нкѝ. Да́ждь гдⷭ҇и и҆ мнѣ̀, недосто́йномꙋ рабꙋ̀ твоемꙋ̀, сп҃се́нїе твоѐ на ло́жи мое́мъ: просвѣтѝ ѹ҆́мъ мо́й свѣ́томъ ра́зꙋма ст҃а́гѡ є҆ђлїа твоегѡ̀, дꙋ́шꙋ любо́вїю крⷭ҇та̀ твоегѡ̀, се́рдце чⷭ҇тото́ю словесѐ твоегѡ̀, тѣ́ло моѐ твое́ю стрⷭ҇тїю безстра́стною: мы́сль мою̀ твои́мъ смире́нїемъ сохранѝ, и҆ воздви́гни мѧ̀ во вре́мѧ подо́бно на твоѐ славосло́вїе. Ꙗ҆́кѡ препросла́вленъ є҆сѝ со безнача́льнымъ твои́мъ ѻ҆ц҃е́мъ, и҆ съ прест҃ы́мъ дх҃омъ, во вѣ́ки. А҆ми́нь.

Мл҃тва г҃, ко прест҃о́мꙋ дх҃ꙋ:
Гдⷭ҇и, цр҃ю̀ нбⷭ҇ный, ѹ҆тѣ́шителю, дш҃е и҆́стины, ѹ҆млⷭ҇рдисѧ, и҆ поми́лꙋй мѧ̀ грѣ́шнаго раба̀ твоего̀, и҆ ѿпꙋсти́ ми недосто́йномꙋ, и҆ простѝ всѧ̑, є҆ли̑ка тѝ согрѣши́хъ дне́сь ꙗ҆́кѡ человѣ́къ, па́че же и҆ не ꙗ҆́кѡ человѣ́къ, но и҆ горѣ́е скота̀: вѡ́льныѧ моѧ̑ грѣхѝ и҆ невѡ́льныѧ, вѣ́дѡмыѧ и҆ невѣ́дѡмыѧ, ꙗ҆̀же ѿ ю҆́ности, и҆ ѿ наꙋ́ки ѕлы̀, и҆ ꙗ҆̀же сꙋ́ть ѿ на́гльства и҆ ѹ҆ны́нїѧ. А҆́ще и҆́менемъ твои́мъ клѧ́хсѧ, и҆лѝ похꙋ́лихъ є҆̀ въ помышле́нїи мое́мъ: и҆лѝ кого̀ ѹ҆кори́хъ, и҆лѝ ѡ҆клевета́хъ кого̀ гнѣ́вомъ мои́мъ, и҆лѝ ѡ҆печа́лихъ: и҆лѝ ѡ҆ че́мъ прогнѣ́вахсѧ, и҆лѝ солга́хъ, и҆лѝ безго́днѡ спа́хъ: и҆лѝ ни́щъ прїи́де ко мнѣ̀, и҆ презрѣ́хъ є҆го̀: и҆лѝ бра́та моего̀ ѡ҆печа́лихъ, и҆лѝ сва́дихъ, и҆лѝ кого̀ ѡ҆сꙋди́хъ: и҆лѝ развелича́хсѧ, и҆лѝ разгордѣ́хсѧ, и҆лѝ разгнѣ́вахсѧ: и҆лѝ стоѧ́щꙋ мѝ на мл҃твѣ, ѹ҆́мъ мо́й ѡ҆ лꙋка́вствїи мі́ра сегѡ̀ подви́жесѧ: и҆лѝ развраще́нїе помы́слихъ, и҆лѝ ѡ҆б꙽ѧдо́хсѧ, и҆лѝ ѡ҆пи́хсѧ, и҆лѝ без꙽ ѹ҆ма̀ смѣѧ́хсѧ: и҆лѝ лꙋка́вое помы́слихъ, и҆лѝ добро́тꙋ чꙋждꙋ́ю ви́дѣвъ, и҆ то́ю ѹ҆ѧ́звленъ бы́хъ се́рдцемъ: и҆лѝ неподѡ́бнаѧ глаго́лахъ, и҆лѝ грѣхꙋ̀ бра́та моегѡ̀ посмѣѧ́хсѧ, моѧ̑ же сꙋ́ть безчи́слєннаѧ согрѣшє́нїѧ: и҆лѝ ѡ҆ мл҃твѣ неради́хъ, и҆лѝ и҆́но что̀ содѣ́ѧхъ лꙋка́вое, не по́мню: та̑ бо всѧ̑ и҆ бѡ́льша си́хъ содѣ́ѧхъ. Поми́лꙋй мѧ̀ тво́рче мо́й влⷣко, ѹ҆ны́лаго и҆ недосто́йнаго раба̀ твоего̀: и҆ ѡ҆ста́ви мѝ, и҆ ѿпꙋстѝ, и҆ прости́ мѧ, ꙗ҆́кѡ бл҃гъ и҆ чл҃вѣколю́бецъ: да съ ми́ромъ лѧ́гꙋ, ѹ҆снꙋ̀ и҆ почі́ю, блꙋ́дный, грѣ́шный и҆ ѻ҆каѧ́нный а҆́зъ: и҆ поклоню́сѧ, и҆ воспою̀, и҆ просла́влю пречⷭ҇тно́е и҆́мѧ твоѐ, со ѻ҆ц҃е́мъ, и҆ є҆диноро́днымъ є҆гѡ̀ сн҃омъ, ны́нѣ и҆ при́снѡ, и҆ во вѣ́ки. А҆ми́нь.

Мл҃тва д҃, ст҃а́гѡ мака́рїа вели́кагѡ:
Что́ ти принесꙋ̀, и҆лѝ что́ ти возда́мъ, великодарови́тый безсме́ртный цр҃ю̀, ще́дре и҆ чл҃вѣколю́бче гдⷭ҇и; ꙗ҆́кѡ лѣнѧ́щасѧ менѐ на твоѐ ѹ҆гожде́нїе, и҆ ничто́же бл҃го сотво́рша, приве́лъ є҆сѝ на коне́цъ мимоше́дшагѡ днѐ сегѡ̀, ѡ҆браще́нїе и҆ сп҃се́нїе дꙋшѝ мое́й стро́ѧ. Млⷭ҇тивъ мѝ бꙋ́ди грѣ́шномꙋ, и҆ ѡ҆бнаже́нномꙋ всѧ́кагѡ дѣ́ла бл҃га: возста́ви па́дшꙋю мою̀ дꙋ́шꙋ, ѡ҆скверни́вшꙋюсѧ въ безмѣ́рныхъ согрѣше́нїихъ, и҆ ѿимѝ ѿ менє̀ ве́сь по́мыслъ лꙋка́вый ви́димагѡ сегѡ̀ житїѧ̀. Простѝ моѧ̑ согрѣшє́нїѧ, є҆ди́не безгрѣ́шне, ꙗ҆́же тѝ согрѣши́хъ въ се́й де́нь, вѣ́дѣнїемъ и҆ невѣ́дѣнїемъ, сло́вомъ, и҆ дѣ́ломъ, и҆ помышле́нїемъ, и҆ всѣ́ми мои́ми чꙋ́вствы. Ты̀ са́мъ покрыва́ѧ сохрани́ мѧ ѿ всѧ́кагѡ сопроти́внагѡ ѡ҆бстоѧ́нїѧ, бжⷭ҇твенною твое́ю вла́стїю, и҆ неизрече́ннымъ чл҃вѣколю́бїемъ, и҆ си́лою. Ѡ҆чи́сти бж҃е, ѡ҆чи́сти мно́жество грѣхѡ́въ мои́хъ. Бл҃говолѝ гдⷭ҇и и҆зба́вити мѧ̀ ѿ сѣ́ти лꙋка́вагѡ, и҆ сп҃сѝ стра́стнꙋю мою̀ дꙋ́шꙋ, и҆ ѡ҆сѣни́ мѧ свѣ́томъ лица̀ твоегѡ̀, є҆гда̀ прїи́деши во сла́вѣ: и҆ неѡсꙋжде́нна ны́нѣ сно́мъ ѹ҆снꙋ́ти сотворѝ, и҆ без꙽ мечта́нїѧ: и҆ несмꙋще́нъ по́мыслъ раба̀ твоегѡ̀ соблюдѝ, и҆ всю̀ сатанинꙋ̀ дѣ́тель ѿженѝ ѿ менє̀: и҆ просвѣти́ ми разꙋ̑мныѧ ѻ҆́чи сердє́чныѧ, да не ѹ҆снꙋ̀ въ сме́рть. И҆ посли́ ми а҆́гг҃ла ми́рна, храни́телѧ и҆ наста́вника дꙋшѝ и҆ тѣ́лꙋ моемꙋ̀, да и҆зба́витъ мѧ̀ ѿ вра̑гъ мои́хъ: да воста́въ со ѻ҆дра̀ моегѡ̀, принесꙋ́ ти бл҃года́рствєнныѧ мольбы̑. Е҆́й гдⷭ҇и, ѹ҆слы́ши мѧ̀ грѣ́шнаго и҆ ѹ҆бо́гаго раба̀ твоего̀, и҆зволе́нїемъ и҆ со́вѣстїю: да́рꙋй мѝ воста́вшꙋ словесє́мъ твои̑мъ поꙋчи́тисѧ, и҆ ѹ҆ны́нїе бѣсо́вское дале́че ѿ менє̀ ѿгна́но бы́ти сотворѝ твои́ми а҆́гг҃лы: да бл҃гословлю̀ и҆́мѧ твоѐ ст҃о́е, и҆ просла́влю, и҆ сла́влю пречⷭ҇тꙋю бцⷣꙋ мр҃і́ю, ю҆́же да́лъ є҆сѝ на́мъ грѣ̑шнымъ застꙋпле́нїе, и҆ прїимѝ сїю̀ молѧ́щꙋюсѧ за ны̀: вѣ́мъ бо, ꙗ҆́кѡ подража́етъ твоѐ чл҃вѣколю́бїе, и҆ молѧ́щисѧ не престае́тъ. Тоѧ̀ застꙋпле́нїемъ, и҆ чⷭ҇тна́гѡ крⷭ҇та̀ зна́менїемъ, и҆ всѣ́хъ ст҃ы́хъ твои́хъ ра́ди, ѹ҆бо́гꙋю дꙋ́шꙋ мою̀ соблюдѝ і҆и҃се хрⷭ҇тѐ бж҃е на́шъ: ꙗ҆́кѡ ст҃ъ є҆сѝ, и҆ препросла́вленъ во вѣ́ки. А҆ми́нь.

Мл҃тва є҃:
Гдⷭ҇и бж҃е на́шъ, є҆́же согрѣши́хъ во днѝ се́мъ сло́вомъ, дѣ́ломъ и҆ помышле́нїемъ, ꙗ҆́кѡ бл҃гъ и҆ чл҃вѣколю́бецъ прости́ ми. Ми́ренъ со́нъ и҆ безмѧте́женъ да́рꙋй мѝ, а҆́гг҃ла твоего̀ храни́телѧ послѝ, покрыва́юща и҆ соблюда́юща мѧ̀ ѿ всѧ́кагѡ ѕла̀: ꙗ҆́кѡ ты̀ є҆сѝ храни́тель дꙋша́мъ и҆ тѣлесє́мъ на́шымъ, и҆ тебѣ̀ сла́вꙋ возсыла́емъ, ѻ҆ц҃ꙋ̀, и҆ сн҃ꙋ, и҆ ст҃о́мꙋ дх҃ꙋ, ны́нѣ и҆ при́снѡ, и҆ во вѣ́ки вѣкѡ́въ. А҆ми́нь.

Мл҃тва ѕ҃:
Гдⷭ҇и бж҃е на́шъ, въ него́же вѣ́ровахомъ, и҆ є҆гѡ́же и҆́мѧ па́че всѧ́кагѡ и҆́мене призыва́емъ: да́ждь на́мъ, ко снꙋ̀ ѿходѧ́щымъ, ѡ҆сла́бꙋ дꙋшѝ и҆ тѣ́лꙋ, и҆ соблюдѝ на́съ ѿ всѧ́кагѡ мечта́нїѧ, и҆ те́мныѧ сла́сти кромѣ̀: ѹ҆ста́ви стремле́нїе страсте́й, ѹ҆гасѝ разжжє́нїѧ воста́нїѧ тѣле́снагѡ. Да́ждь на́мъ цѣломꙋ́дреннѣ пожи́ти дѣ́лы и҆ словесы̀: да добродѣ́тельное жи́тельство воспрїе́млюще, ѡ҆бѣтова́нныхъ не ѿпаде́мъ бл҃ги́хъ твои́хъ, ꙗ҆́кѡ бл҃гослове́нъ є҆сѝ во вѣ́ки. А҆ми́нь.

Мл҃тва з҃, ст҃а́гѡ і҆ѡа́нна златоꙋ́стагѡ: мл҃твы молє́бныѧ, число́мъ два́десѧтимъ четы́ремъ часѡ́мъ, дневны̑мъ и҆ нощны̑мъ.
[№] Гдⷭ҇и, не лишѝ менє̀ нбⷭ҇ныхъ твои́хъ бла̑гъ. [в҃] Гдⷭ҇и, и҆зба́ви мѧ̀ вѣ́чныхъ мꙋ́къ. [г҃] Гдⷭ҇и, ѹ҆мо́мъ ли и҆лѝ помышле́нїемъ, сло́вомъ и҆лѝ дѣ́ломъ согрѣши́хъ, прости́ мѧ. [д҃] Гдⷭ҇и, и҆зба́ви мѧ̀ всѧ́кагѡ невѣ́дѣнїѧ, и҆ забве́нїѧ, и҆ малодꙋ́шїѧ, и҆ ѡ҆камене́ннагѡ нечꙋ́вствїѧ. [є҃] Гдⷭ҇и, и҆зба́ви мѧ̀ ѿ всѧ́кагѡ и҆скꙋше́нїѧ. [ѕ҃] Гдⷭ҇и, просвѣтѝ моѐ се́рдце, є҆́же помрачѝ лꙋка́вое похотѣ́нїе. [з҃] Гдⷭ҇и, а҆́зъ ꙗ҆́кѡ человѣ́къ согрѣши́хъ, ты́ же ꙗ҆́кѡ бг҃ъ ще́дръ, поми́лꙋй мѧ̀, ви́дѧ не́мощь дꙋшѝ моеѧ̀. [и҃] Гдⷭ҇и, послѝ бл҃года́ть твою̀ въ по́мощь мнѣ̀, да просла́влю и҆́мѧ твоѐ ст҃о́е. [ѳ҃] Гдⷭ҇и і҆и҃се хрⷭ҇тѐ, напиши́ мѧ раба̀ твоего̀ въ кни́зѣ живо́тнѣй, и҆ да́рꙋй мѝ коне́цъ бл҃гі́й. [‹] Гдⷭ҇и бж҃е мо́й, а҆́ще и҆ ничто́же бл҃го сотвори́хъ пред꙽ тобо́ю, но да́ждь мѝ по бл҃года́ти твое́й, положи́ти нача́ло бл҃го́е. [№i] Гдⷭ҇и, ѡ҆кропѝ въ се́рдцѣ мое́мъ ро́сꙋ бл҃года́ти твоеѧ̀. [в҃i] Гдⷭ҇и нб҃сѐ и҆ землѝ, помѧни́ мѧ грѣ́шнаго раба̀ твоего̀, стꙋ́днаго и҆ нечи́стаго, во црⷭ҇твїи твое́мъ. А҆ми́нь.

[№] Гдⷭ҇и, въ покаѧ́нїи прїими́ мѧ. [в҃] Гдⷭ҇и, не ѡ҆ста́ви менє̀. [г҃] Гдⷭ҇и, не введѝ менє̀ въ напа́сть. [д҃] Гдⷭ҇и, да́ждь мѝ мы́сль бл҃гꙋ. [є҃] Гдⷭ҇и, да́ждь мѝ сле́зы, и҆ па́мѧть сме́ртнꙋю, и҆ ѹ҆миле́нїе. [ѕ҃] Гдⷭ҇и, да́ждь мѝ по́мыслъ и҆сповѣ́данїѧ грѣхѡ́въ мои́хъ. [з҃] Гдⷭ҇и, да́ждь мѝ смире́нїе, цѣломꙋ́дрїе и҆ послꙋша́нїе. [и҃] Гдⷭ҇и, да́ждь мѝ терпѣ́нїе, великодꙋ́шїе и҆ кро́тость. [ѳ҃] Гдⷭ҇и, вселѝ въ мѧ̀ ко́рень бл҃ги́хъ, стра́хъ тво́й въ се́рдце моѐ. [‹] Гдⷭ҇и, сподо́би мѧ̀ люби́ти тѧ̀ ѿ всеѧ̀ дꙋшѝ моеѧ̀ и҆ помышле́нїѧ, и҆ твори́ти во все́мъ во́лю твою̀. [№i] Гдⷭ҇и, покры́й мѧ̀ ѿ человѣ̑къ нѣ́которыхъ, и҆ бѣсѡ́въ, и҆ страсте́й, и҆ ѿ всѧ́кїѧ и҆ны́ѧ неподо́бныѧ ве́щи. [в҃i] Гдⷭ҇и, вѣ́си, ꙗ҆́кѡ твори́ши, ꙗ҆́коже ты̀ во́лиши, да бꙋ́детъ во́лѧ твоѧ̀ и҆ во мнѣ̀ грѣ́шнѣмъ: ꙗ҆́кѡ бл҃гослове́нъ є҆сѝ во вѣ́ки. А҆ми́нь.

Мл҃тва и҃, ко гдⷭ҇ꙋ на́шемꙋ і҆и҃сꙋ хрⷭ҇тꙋ̀:
Гдⷭ҇и і҆и҃се хрⷭ҇тѐ сн҃е бж҃їй, ра́ди чⷭ҇тнѣ́йшїѧ мт҃ре твоеѧ̀, и҆ безпло́тныхъ твои́хъ а҆́гг҃лъ, прⷪ҇ро́ка же и҆ п®те́чи и҆ крⷭ҇ти́телѧ твоегѡ̀, бг҃оглаго́ливыхъ же а҆пⷭ҇лъ, свѣ́тлыхъ и҆ добропобѣ́дныхъ мч҃нкъ, прпⷣбныхъ и҆ бг҃оно́сныхъ ѻ҆тє́цъ, и҆ всѣ́хъ ст҃ы́хъ мл҃твами, и҆зба́ви мѧ̀ настоѧ́щагѡ ѡ҆бстоѧ́нїѧ бѣсо́вскагѡ. Е҆́й, гдⷭ҇и мо́й и҆ тво́рче, не хотѧ́й сме́рти грѣ́шнагѡ, но ꙗ҆́коже ѡ҆брати́тисѧ и҆ жи́вꙋ бы́ти є҆мꙋ̀, да́ждь и҆ мнѣ̀ ѡ҆браще́нїе ѻ҆каѧ́нномꙋ и҆ недосто́йномꙋ: и҆зми́ мѧ ѿ ѹ҆́стъ па́гꙋбнагѡ ѕмі́ѧ, зїѧ́ющагѡ пожре́ти мѧ̀, и҆ свестѝ во а҆́дъ жи́ва. Е҆́й, гдⷭ҇и мо́й, ѹ҆тѣше́нїе моѐ, и҆́же менє̀ ра́ди ѻ҆каѧ́ннагѡ въ тлѣ́ннꙋю пло́ть ѡ҆болкі́йсѧ, и҆сто́ргни мѧ̀ ѿ ѻ҆каѧ́нства, и҆ ѹ҆тѣше́нїе пода́ждь дꙋшѝ мое́й ѻ҆каѧ́ннѣй. Всадѝ въ се́рдце моѐ твори́ти твоѧ̑ повелѣ̑нїѧ, и҆ ѡ҆ста́вити лꙋка̑ваѧ дѣѧ̑нїѧ, и҆ полꙋчи́ти бл҃жє́нства твоѧ̑: на тѧ́ бо гдⷭ҇и ѹ҆пова́хъ, сп҃си́ мѧ.

Мл҃тва ѳ҃, ко прест҃ѣ́й бцⷣѣ,
ст҃а́гѡ петра̀ стꙋді́йскагѡ:
Къ тебѣ̀ пречтⷭ҇ѣй бж҃їей мт҃ри а҆́зъ ѻ҆каѧ́нный припа́даѧ молю́сѧ: вѣ́си, цр҃и́це, ꙗ҆́кѡ безпреста́ни согрѣша́ю и҆ прогнѣвлѧ́ю сн҃а твоего̀ и҆ бг҃а моего̀, и҆ мно́гажды а҆́ще ка́юсѧ, ло́жъ пред꙽ бг҃омъ ѡ҆брѣта́юсѧ, и҆ ка́юсѧ трепе́щѧ: не ѹ҆же́ли гдⷭ҇ь порази́тъ мѧ̀, и҆ по часѣ̀ па́ки та̑ѧжде творю̀. Вѣ́дꙋщи сїѧ̑, влⷣчце моѧ̀ гпⷭ҇жѐ бцⷣе, молю̀, да поми́лꙋеши, да ѹ҆крѣпи́ши, и҆ бл҃га̑ѧ твори́ти да пода́си мѝ. Вѣ́си бо, влⷣчце моѧ̀ бцⷣе, ꙗ҆́кѡ ѿню́дъ и҆́мамъ въ не́нависти ѕла̑ѧ моѧ̑ дѣла̀, и҆ все́ю мы́слїю люблю̀ зако́нъ бг҃а моегѡ̀: но не вѣ́мъ, гпⷭ҇жѐ пречтⷭ҇аѧ, ѿкꙋ́дꙋ ꙗ҆̀же ненави́ждꙋ, та̑ и҆ люблю̀, а҆ бл҃га̑ѧ престꙋпа́ю. Не попꙋща́й, пречтⷭ҇аѧ, во́ли мое́й соверша́тисѧ, не ѹ҆го́дна бо є҆́сть: но да бꙋ́детъ во́лѧ сн҃а твоегѡ̀ и҆ бг҃а моегѡ̀: да мѧ̀ сп҃се́тъ, и҆ вразꙋми́тъ, и҆ пода́стъ бл҃года́ть ст҃а́гѡ дх҃а, да бы́хъ а҆́зъ ѿсе́лѣ преста́лъ сквернодѣ́йства, и҆ про́чее пожи́лъ бы́хъ въ повелѣ́нїи сн҃а твоегѡ̀, є҆мꙋ́же подоба́етъ всѧ́каѧ сла́ва, че́сть и҆ держа́ва, со безнача́льнымъ є҆гѡ̀ ѻ҆ц҃е́мъ, и҆ прест҃ы́мъ, и҆ бл҃ги́мъ, и҆ животворѧ́щимъ є҆гѡ̀ дх҃омъ, ны́нѣ и҆ при́снѡ, и҆ во вѣ́ки вѣкѡ́въ. А҆ми́нь.

Мл҃тва ‹, ко прест҃ѣ́й бцⷣѣ:
Бл҃га́гѡ цр҃ѧ̀ бл҃га́ѧ мт҃и, пречтⷭ҇аѧ и҆ бл҃гослове́ннаѧ бцⷣе мр҃і́е, млⷭ҇ть сн҃а твоегѡ̀ и҆ бг҃а на́шегѡ и҆зле́й на стра́стнꙋю мою̀ дꙋ́шꙋ, и҆ твои́ми мл҃твами наста́ви мѧ̀ на дѣѧ̑нїѧ бл҃га̑ѧ: да про́чее вре́мѧ живота̀ моегѡ̀ без꙽ поро́ка прейдꙋ̀, и҆ тобо́ю ра́й да ѡ҆брѧ́щꙋ, бцⷣе дв҃о, є҆ди́на чтⷭ҇аѧ и҆ бл҃гослове́ннаѧ.

Мл҃тва №i, ко ст҃о́мꙋ а҆́гг҃лꙋ храни́телю:
А҆́гг҃ле хрⷭ҇то́въ, храни́телю мо́й ст҃ы́й, и҆ покрови́телю дꙋшѝ и҆ тѣ́ла моегѡ̀, всѧ̑ мѝ простѝ, є҆ли̑ка согрѣши́хъ во дне́шнїй де́нь, и҆ ѿ всѧ́кагѡ лꙋка́вствїѧ проти́внагѡ мѝ врага̀ и҆зба́ви мѧ̀, да ни въ ко́емже грѣсѣ̀ прогнѣ́ваю бг҃а моего̀: но молѝ за мѧ̀ грѣ́шнаго и҆ недосто́йнаго раба̀, ꙗ҆́кѡ да досто́йна мѧ̀ пока́жеши бл҃гости и҆ млⷭ҇ти всест҃ы́ѧ трⷪ҇цы, и҆ мт҃ре гдⷭ҇а моегѡ̀ і҆и҃са хрⷭ҇та̀, и҆ всѣ́хъ ст҃ы́хъ. А҆ми́нь.

Конда́къ бцⷣѣ, гла́съ и҃:
Взбра́нной воево́дѣ побѣди́тєльнаѧ, ꙗ҆́кѡ и҆зба́вльшесѧ ѿ ѕлы́хъ, бл҃года́рствєннаѧ восписꙋ́емъ тѝ рабѝ твоѝ, бцⷣе: но ꙗ҆́кѡ и҆мꙋ́щаѧ держа́вꙋ непобѣди́мꙋю, ѿ всѧ́кихъ на́съ бѣ́дъ свободѝ, да зове́мъ тѝ: ра́дꙋйсѧ, невѣ́сто неневѣ́стнаѧ.

Пресла́внаѧ приснодв҃о, мт҃и хрⷭ҇та̀ бг҃а, принесѝ на́шꙋ мл҃твꙋ сн҃ꙋ твоемꙋ̀ и҆ бг҃ꙋ на́шемꙋ, да сп҃се́тъ тобо́ю дꙋ́шы на́шѧ.

Всѐ ѹ҆пова́нїе моѐ на тѧ̀ возлага́ю, мт҃и бж҃їѧ, сохрани́ мѧ под꙽ кро́вомъ твои́мъ.

Бцⷣе дв҃о, не пре́зри менє̀ грѣ́шнагѡ, тре́бꙋюща твоеѧ̀ по́мощи и҆ твоегѡ̀ застꙋпле́нїѧ: на тѧ́ бо ѹ҆пова̀ дꙋша̀ моѧ̀, и҆ поми́лꙋй мѧ̀.

Мл҃тва ст҃а́гѡ і҆ѡанні́кїа:
Ѹ҆пова́нїе моѐ ѻ҆ц҃ъ, прибѣ́жище моѐ сн҃ъ, покро́въ мо́й дх҃ъ ст҃ы́й: трⷪ҇це ст҃а́ѧ, сла́ва тебѣ̀.

Ѡ҆кончанїе мл҃твъ:
Досто́йно є҆́сть ꙗ҆́кѡ вои́стиннꙋ, бл҃жи́ти тѧ̀ бцⷣꙋ, приснобл҃же́ннꙋю, и҆ пренепоро́чнꙋю, и҆ мт҃рь бг҃а на́шегѡ. ЧCтнѣ́йшꙋю херꙋві̑мъ, и҆ сла́внѣйшꙋю без꙽ сравне́нїѧ серафі̑мъ, без꙽ и҆стлѣ́нїѧ бг҃а сло́ва ро́ждшꙋю, сꙋ́щꙋю бцⷣꙋ тѧ̀ велича́емъ.

Сла́ва ѻ҆ц҃ꙋ̀, и҆ сн҃ꙋ, и҆ ст҃о́мꙋ дх҃ꙋ, и҆ ны́нѣ и҆ при́снѡ, и҆ во вѣ́ки вѣкѡ́въ. А҆ми́нь.

Гдⷭ҇и, поми́лꙋй. [Три́жды.]

Гдⷭ҇и і҆и҃се хрⷭ҇тѐ сн҃е бж҃їй, мл҃твъ ра́ди пречⷭ҇тыѧ твоеѧ̀ мт҃ре, прпⷣбныхъ и҆ бг҃оно́сныхъ ѻ҆тє́цъ на́шихъ, и҆ всѣ́хъ ст҃ы́хъ, поми́лꙋй на́съ. А҆ми́нь.

Та́же, вмѣ́стѡ проще́нїѧ:
Ѡ҆сла́би, ѡ҆ста́ви, простѝ бж҃е прегрѣшє́нїѧ на̑ша, вѡ́льнаѧ и҆ невѡ́льнаѧ, ꙗ҆̀же въ сло́вѣ и҆ въ дѣ́лѣ, ꙗ҆̀же въ вѣ́дѣнїи и҆ не въ вѣ́дѣнїи, ꙗ҆̀же во днѝ и҆ въ нощѝ, ꙗ҆̀же во ѹ҆мѣ̀ и҆ въ помышле́нїи: всѧ̑ на́мъ простѝ, ꙗ҆́кѡ бл҃гъ и҆ чл҃вѣколю́бецъ.

Мл҃тва:
Ненави́дѧщихъ и҆ ѡ҆би́дѧщихъ на́съ простѝ, гдⷭ҇и чл҃вѣколю́бче. Бл҃готворѧ́щымъ бл҃госотворѝ. Бра́тїѧмъ и҆ сро́дникѡмъ на́шымъ да́рꙋй ꙗ҆̀же ко сп҃се́нїю прошє́нїѧ, и҆ жи́знь вѣ́чнꙋю. Въ не́мощехъ сꙋ́щыѧ посѣтѝ, и҆ и҆сцѣле́нїе да́рꙋй. Я%же въ мо́ри ѹ҆пра́ви. Пꙋтеше́ствꙋющымъ спꙋтеше́ствꙋй. Хрⷭ҇тїа́нѡмъ спобо́рствꙋй. Слꙋжа́щымъ и҆ ми́лꙋющымъ на́съ грѣхѡ́въ ѡ҆ставле́нїе да́рꙋй. Заповѣ́давшихъ на́мъ недостѡ́йнымъ моли́тисѧ ѡ҆ ни́хъ, поми́лꙋй по вели́цѣй твое́й млⷭ҇ти. Помѧнѝ гдⷭ҇и, пре́жде ѹ҆со́пшихъ ѻ҆тє́цъ и҆ бра́тїй на́шихъ, и҆ ѹ҆поко́й и҆̀хъ, и҆дѣ́же присѣща́етъ свѣ́тъ лица̀ твоегѡ̀. Помѧнѝ гдⷭ҇и, бра́тїй на́шихъ плѣне́нныхъ, и҆ и҆зба́ви ѧ҆̀ ѿ всѧ́кагѡ ѡ҆бстоѧ́нїѧ. Помѧнѝ гдⷭ҇и, плодоносѧ́щихъ и҆ добродѣ́лающихъ во ст҃ы́хъ твои́хъ цр҃квахъ, и҆ да́ждь и҆̀мъ ꙗ҆́же ко сп҃се́нїю прошє́нїѧ, и҆ жи́знь вѣ́чнꙋю. Помѧнѝ гдⷭ҇и и҆ на́съ, смире́нныхъ, и҆ грѣ́шныхъ, и҆ недосто́йныхъ ра̑бъ твои́хъ, и҆ просвѣтѝ на́шъ ѹ҆́мъ свѣ́томъ ра́зꙋма твоегѡ̀, и҆ наста́ви на́съ на стезю̀ за́повѣдей твои́хъ: мл҃твами пречⷭ҇тыѧ влⷣчцы на́шеѧ бцⷣы и҆ приснодв҃ы мр҃і́и, и҆ всѣ́хъ твои́хъ ст҃ы́хъ, ꙗ҆́кѡ бл҃гослове́нъ є҆сѝ во вѣ́ки вѣкѡ́въ. А҆ми́нь.

И҆сповѣ́данїе грѣхѡ́въ повседне́вное:
И҆сповѣ́даю тебѣ̀ гдⷭ҇ꙋ бг҃ꙋ моемꙋ̀ и҆ творцꙋ̀, во ст҃ѣ́й трⷪ҇цѣ є҆ди́номꙋ, сла́вимомꙋ и҆ покланѧ́емомꙋ, ѻ҆ц҃ꙋ̀, и҆ сн҃ꙋ, и҆ ст҃о́мꙋ дх҃ꙋ, всѧ̑ моѧ̑ грѣхѝ, ꙗ҆̀же содѣ́ѧхъ во всѧ̑ дни̑ живота̀ моегѡ̀, и҆ на всѧ́кїй ча́съ, и҆ въ настоѧ́щее вре́мѧ, и҆ въ преше́дшыѧ дни̑ и҆ нѡ́щи: дѣ́ломъ, сло́вомъ, помышле́нїемъ, ѡ҆б꙽ѧде́нїемъ, пїѧ́нствомъ, тайноѧде́нїемъ, праздносло́вїемъ, ѹ҆ны́нїемъ, лѣ́ностїю, прекосло́вїемъ, непослꙋша́нїемъ, ѡ҆клевета́нїемъ, ѡ҆сꙋжде́нїемъ, небреже́нїемъ, самолю́бїемъ, многостѧжа́нїемъ, хище́нїемъ, неправдоглаго́ланїемъ, скверноприбы́тчествомъ, мшелои́мствомъ, ревнова́нїемъ, за́вистїю, гнѣ́вомъ, памѧтоѕло́бїемъ, не́навистїю, лихои́мствомъ, и҆ всѣ́ми мои́ми чꙋ́вствы: зрѣ́нїемъ, слꙋ́хомъ, ѡ҆бонѧ́нїемъ, вкꙋ́сомъ, ѡ҆сѧза́нїемъ, и҆ про́чими мои́ми грѣхи̑, дꙋше́вными вкꙋ́пѣ и҆ тѣле́сными, и҆́миже тебѐ бг҃а моего̀ и҆ творца̀ прогнѣ́вахъ, и҆ бли́жнѧго моего̀ ѡ҆непра́вдовахъ. Ѡ҆ си́хъ жалѣ́ѧ, ви́нна себѐ тебѣ̀, бг҃ꙋ моемꙋ̀, представлѧ́ю, и҆ и҆мѣ́ю во́лю ка́ѧтисѧ: то́чїю, гдⷭ҇и бж҃е мо́й, помози́ ми, со слеза́ми смире́ннѡ молю́ тѧ: прешє́дшаѧ же согрѣшє́нїѧ моѧ̑ млⷭ҇рдїемъ твои́мъ прости́ ми, и҆ разрѣшѝ ѿ всѣ́хъ си́хъ, ꙗ҆̀же и҆зглаго́лахъ пред꙽ тобо́ю, ꙗ҆́кѡ бл҃гъ и҆ чл҃вѣколю́бецъ.

Мл҃тва ст҃а́гѡ і҆ѡа́нна дамаски́на,
ю҆́же ѹ҆казꙋ́ѧ на ѻ҆́дръ тво́й, глаго́ли:
Влⷣко чл҃вѣколю́бче, не ѹ҆же́ли мнѣ̀ ѻ҆́дръ се́й гро́бъ бꙋ́детъ; и҆лѝ є҆щѐ ѻ҆каѧ́ннꙋю мою̀ дꙋ́шꙋ просвѣти́ши дне́мъ; Се́ ми гро́бъ предлежи́тъ, се́ ми сме́рть предстои́тъ. Сꙋда̀ твоегѡ̀ гдⷭ҇и бою́сѧ, и҆ мꙋ́ки безконе́чныѧ, ѕло́е же творѧ̀ не престаю̀: тебѐ гдⷭ҇а бг҃а моего̀ всегда̀ прогнѣвлѧ́ю, и҆ пречⷭ҇тꙋю твою̀ мт҃рь, и҆ всѧ̑ нбⷭ҇ныѧ си̑лы, и҆ ст҃а́го а҆́гг҃ла храни́телѧ моего̀. Вѣ́мъ ѹ҆̀бо гдⷭ҇и, ꙗ҆́кѡ недосто́инъ є҆́смь чл҃вѣколю́бїѧ твоегѡ̀, но досто́инъ є҆́смь всѧ́кагѡ ѡ҆сꙋжде́нїѧ и҆ мꙋ́ки. Но гдⷭ҇и, и҆лѝ хощꙋ̀, и҆лѝ не хощꙋ̀, сп҃си́ мѧ. А҆́ще бо пра́ведника сп҃се́ши, ничто́же ве́лїе, и҆ а҆́ще чи́стаго поми́лꙋеши, ничто́же ди́вно: досто́йни бо сꙋ́ть млⷭ҇ти твоеѧ̀. Но на мнѣ̀ грѣ́шнѣмъ ѹ҆дивѝ млⷭ҇ть твою̀: ѡ҆ се́мъ ꙗ҆вѝ чл҃вѣколю́бїе твоѐ, да не ѡ҆долѣ́етъ моѧ̀ ѕло́ба твое́й неизглаго́ланнѣй бл҃гости и҆ млⷭ҇рдїю: и҆ ꙗ҆́коже хо́щеши, ѹ҆стро́й ѡ҆ мнѣ̀ ве́щь.

И҆ та́кѡ хотѧ́й возлещѝ на посте́лю,
глаго́ли тропарѝ сїѧ̑, гла́съ в҃:
Просвѣтѝ ѻ҆́чи моѝ хрⷭ҇тѐ бж҃е: да не когда̀ ѹ҆снꙋ̀ въ сме́рть, да не когда̀ рече́тъ вра́гъ мо́й: ѹ҆крѣпи́хсѧ на него̀.

Сла́ва ѻ҆ц҃ꙋ̀, и҆ сн҃ꙋ, и҆ ст҃о́мꙋ дх҃ꙋ.

Застꙋ́пникъ дꙋшѝ моеѧ̀ бꙋ́ди бж҃е, ꙗ҆́кѡ посредѣ̀ хождꙋ̀ сѣте́й мно́гихъ: и҆зба́ви мѧ̀ ѿ ни́хъ, и҆ сп҃си́ мѧ бл҃же, ꙗ҆́кѡ чл҃вѣколю́бецъ.

И҆ ны́нѣ и҆ при́снѡ, и҆ во вѣ́ки вѣкѡ́въ. А҆ми́нь.

Пресла́внꙋю бж҃їю мт҃рь, и҆ ст҃ы́хъ а҆́гг҃лъ ст҃ѣ́йшꙋю, немо́лчнѡ воспои́мъ се́рдцемъ и҆ ѹ҆сты̀, бцⷣꙋ сїю̀ и҆сповѣ́дающе, ꙗ҆́кѡ вои́стиннꙋ ро́ждшꙋю на́мъ бг҃а воплоще́нна, и҆ молѧ́щꙋюсѧ непреста́ннѡ ѡ҆ дꙋша́хъ на́шихъ.

Та́же цѣлꙋ́й крⷭ҇тъ тво́й, и҆ прекрестѝ крⷭ҇то́мъ мѣ́сто твоѐ ѿ главы̀ и҆ до но́гъ, та́кожде и҆ ѿ всѣ́хъ стра́нъ, глаго́лѧ:

Мл҃тва чⷭ҇тно́мꙋ крⷭ҇тꙋ̀:
Да воскрⷭ҇нетъ бг҃ъ, и҆ расточа́тсѧ вразѝ є҆гѡ̀, и҆ да бѣжа́тъ ѿ лица̀ є҆гѡ̀ ненави́дѧщїи є҆го̀. Ꙗ҆́кѡ и҆счеза́етъ ды́мъ, да и҆сче́знꙋтъ, ꙗ҆́кѡ та́етъ во́скъ ѿ лица̀ ѻ҆гнѧ̀, та́кѡ да поги́бнꙋтъ бѣ́си ѿ лица̀ лю́бѧщихъ бг҃а, и҆ зна́менꙋющихсѧ крⷭ҇тнымъ зна́менїемъ, и҆ въ весе́лїи глаго́лющихъ: ра́дꙋйсѧ, пречⷭ҇тны́й и҆ животворѧ́щїй крⷭ҇те гдⷭ҇ень, прогонѧ́ѧй бѣ́сы си́лою на тебѣ̀ пропѧ́тагѡ гдⷭ҇а на́шегѡ і҆и҃са хрⷭ҇та̀, во а҆́дъ сше́дшагѡ, и҆ попра́вшагѡ си́лꙋ дїа́волю, и҆ дарова́вшагѡ на́мъ тебѐ, крⷭ҇тъ сво́й чⷭ҇тны́й, на прогна́нїе всѧ́кагѡ сꙋпоста́та. Q пречⷭ҇тны́й и҆ животворѧ́щїй крⷭ҇те гдⷭ҇ень! помога́й мѝ со ст҃о́ю гпⷭ҇же́ю дв҃ою бцⷣею, и҆ со всѣ́ми ст҃ы́ми во вѣ́ки. А҆ми́нь.

И҆ли кра́ткѡ:
Ѡ҆гради́ мѧ гдⷭ҇и си́лою чⷭ҇тна́гѡ и҆ животворѧ́щагѡ твоегѡ̀ крⷭ҇та̀, и҆ сохрани́ мѧ ѿ всѧ́кагѡ ѕла̀.

Е\#гда̀ ѿхо́диши ко снꙋ̀, глаго́ли:
Въ рꙋ́цѣ твоѝ, гдⷭ҇и і҆и҃се хрⷭ҇тѐ бж҃е мо́й, предаю̀ дꙋ́хъ мо́й: ты́ же мѧ̀ бл҃гословѝ, ты̀ мѧ̀ поми́лꙋй, и҆ живо́тъ вѣ́чный да́рꙋй мѝ. А҆ми́нь.

\end{document}
